\section{Introduction}

Federated learning (FL) trains a shared model across many clients without
centralizing raw data \citep{mcmahan2017}. In practice, however, participating
devices are rarely homogeneous. A laptop-class client may upload a
full-precision update, while an embedded device or phone may need INT8, INT4,
or even INT2 quantization to fit its communication and memory budget. Once
those updates are averaged together, quantization error no longer behaves like
benign noise; it becomes a structured bias in the aggregated direction.

This project studies a simple question: can a server learn to reverse some of
that bias if a small subset of clients still provide clean, full-precision
updates? We use those full-precision clients as \emph{anchors}. Their updates
provide round-level supervision for a server-side CorrectionNet that maps a
quantized delta toward a more full-precision-like delta before aggregation.

That question sits at the intersection of two existing FL lines of work.
Heterogeneity-aware optimizers such as FedProx, SCAFFOLD, and FedNova study how
client drift and uneven local progress affect aggregation
\citep{li2020fedprox,karimireddy2020scaffold,wang2020fednova}. Mixed-precision
and bit-width-aware methods such as ProWD, FedShift, and recent client-specific
quantization schemes instead focus on what happens when clients cannot all
communicate at the same precision
\citep{yoon2022prowd,seo2024fedshift,chen2024fedmpq}.

The current implementation goes beyond a naive dequantizer in several ways.
Instead of training on only the current round, it keeps a replay buffer over the
most recent rounds. Instead of forcing every low-bit client toward the same
anchor mean, it can assign each constrained client its most similar anchor by
cosine similarity. It also supports an optional distillation term computed on a
server validation set and blends corrected deltas with the original quantized
delta to preserve client-specific direction under non-IID data.

Our main contributions are:
\begin{itemize}[leftmargin=*]
  \item A template-based writeup of the implemented anchor-guided correction
  pipeline, including warmup, replay buffering, smart matching, normalization,
  and optional similarity-aware aggregation.
  \item A reusable architecture diagram that makes the per-round data flow and
  supervision path explicit.
  \item An empirical summary showing that the method is most useful in the
  severe quantization regime, especially INT2, while still improving over the
  simplified ProWD baseline in every reported configuration and matching or
  beating FedShift on the shared settings, with the clearest gains at early
  INT2, final INT4, and non-IID INT8.
\end{itemize}
